\subsection{Overview}

A large body of plot-, field-, watershed-, and regional-scale evidence converges on a consistent picture of how agronomic management governs nitrate-nitrogen (NO$_3^-$-N) concentrations and loads in agricultural subsurface (tile) drainage. Across the U.S. Corn Belt and comparable humid systems in Canada, Europe, and Asia, drainage nitrate loss is driven jointly by nitrogen (N) input, crop and rotation choice, the amount and timing of drainage flow, and soil processing of N, so that no single practice is universally decisive but several act in predictable directions \parencite{Dinnes2002}. The strongest and most repeatable levers are the amount of N applied relative to crop demand, the presence of a living or perennial crop to take up residual N during the drainage season, and the fraction of the year in which fields drain. Under conventional maize and maize--soybean systems, flow-weighted NO$_3^-$-N concentrations commonly exceed the 10~mg~L$^{-1}$ drinking-water standard: more than 70\% of site-years in one twelve-state synthesis exceeded it \parencite{Liu2021}, and corn, corn--soybean, and soybean site-years averaged 14.0, 13.5, and 12.1~mg~L$^{-1}$ across a continental database \parencite{Hertzberger2019}.

Several headline findings are well supported. First, drainage nitrate concentration and load increase with N application rate, while yield plateaus near the economic optimum, so that reducing rate toward the agronomic/economic optimum lowers loss with modest yield penalty \parencite{Jaynes2001,Helmers2012b,Zhao2016,Hertzberger2019}. Second, cover crops---especially winter cereal rye---and perennial or diversified rotations are among the most consistent and largest in-field reducers of drainage nitrate, frequently cutting loads by one-third to more than one-half \parencite{Ruffatti2019,OBrien2022,PrezaFontes2021,Daigh2015b,Gentry2025}. Third, shifting N from fall to spring or split/side-dress timing, and using soil-test-based rate programs, reduces loss but with smaller and less consistent effects than rate or cropping-system changes \parencite{Randall2003,Dinnes2002,Jaynes2004}. Fourth, tillage exerts weak and inconsistent control over nitrate loss relative to precipitation, cropping system, and N rate \parencite{RANDALL1995,Weed1996,Randall2001}. Fifth, hydrology---the amount and seasonal timing of drainage, itself governed by precipitation---is the dominant control on interannual variability, with most annual loss occurring in spring \parencite{Randall2003,Bakhsh2010,RANDALL1995}. Finally, mitigation practices can shift N among loss pathways (leaching, N$_2$O, NH$_3$), so evaluating them on a common environmental basis is increasingly emphasized \parencite{PrezaFontes2023,Menegaz2025}. The subsections below synthesize this evidence by management lever.

\subsection{Nitrogen Application Rate and Residual Soil Nitrate}

Drainage nitrate loss rises with N application rate, and this is the most consistently documented in-field relationship. In a 22-ha Iowa production field, yearly nitrate mass loss was significantly greater for a high-N treatment (48~kg~N~ha$^{-1}$) than for medium (35) and low (29) treatments, and concentrations ranged 5--45~mg~N~L$^{-1}$, exceeding the 10~mg~N~L$^{-1}$ maximum contaminant level except under low N during soybean years \parencite{Jaynes2001}. In north-central Iowa, flow-weighted NO$_3^-$-N concentrations increased with rate in both continuous corn and corn--soybean systems, four-year average losses ranged 22--68~kg~N~ha$^{-1}$, and concentrations at recommended rates commonly exceeded 10~mg~L$^{-1}$ \parencite{Helmers2012b}. Early Iowa work found losses rose with rate (50--200~kg~ha$^{-1}$) but not linearly, identifying a single 150~kg~ha$^{-1}$ application as a favorable compromise between loss and yield \parencite{Azevedo1997}. Increasing N above the agronomic optimum yielded diminishing returns: at 168--252~kg~N~ha$^{-1}$ in Illinois, maize yield differed little (15.3, 14.5, 15.1, 14.9~Mg~ha$^{-1}$) while flow-weighted concentration was highest at the lowest rate in one comparison and cumulative loads (14.4--22.8~kg~ha$^{-1}$) were statistically similar \parencite{Li2023b}. A ten-percent rate reduction cut Illinois corn yield about 10\% but did not significantly reduce tile nitrate loss in one study, underscoring that rate response varies with site and hydrology \parencite{Gentry2024}.

Regional and database syntheses formalize these patterns. Across the United States and Canada, maize yield followed a quadratic response peaking at 9.3~Mg~ha$^{-1}$ near 193~kg~N~ha$^{-1}$, drainage N losses increased exponentially with rate, and yield-scaled N loss rose sharply once N surplus exceeded about 23~kg~N~ha$^{-1}$---identifying zero N surplus as a practical target \parencite{Zhao2016}. In the MANAGE drain-concentration database, N rate was the strongest predictor of annual NO$_3^-$-N concentration in corn site-years ($R^2=0.59$), with the $>224$~kg~ha$^{-1}$ category highest (19.2~mg~L$^{-1}$) and no yield gain above 150--224~kg~ha$^{-1}$ \parencite{Hertzberger2019}. A twelve-state analysis confirmed that higher rates raised yield, load, and concentration, that reducing rate mitigated loss but reduced yield, and that the most cost-effective high-rate scenarios also produced the highest losses \parencite{Liu2021}. A meta-analysis of ~1,000 site-years found dissolved N load increased with rate in continuous corn but not corn--soybean \parencite{Christianson2015b}, and a 62-study synthesis reported each unit rise in fertilizer-N rate produced 0.12~kg~NO$_3^-$-N~ha$^{-1}$ loss, with rate explaining 27\% of nitrate-loss variability \parencite{Chatterjee2020b}. Historical measurements likewise scaled loads with rate---e.g., 41, 55, 93, and 180~kg~ha$^{-1}$ at 20--448~kg~ha$^{-1}$ applied \parencite{Ritter2015}.

Residual soil nitrate (RSN) links rate to leaching risk. Across eight states, RSN was minimally affected by rates below the economic optimum N rate (EONR) and increased sharply only above it, following a plateau-quadratic form; at EONR, RSN was lower for at-plant than split applications (42 vs.\ 62~kg~N~ha$^{-1}$), and RSN did not substantially rise until rates exceeded EONR by more than 30~kg~N~ha$^{-1}$ \parencite{Bandura2025}. Long-term Minnesota data showed continuous corn accumulated up to 225~kg~ha$^{-1}$ RSN in the profile during dry years, then flushed to drainage in subsequent wet years \parencite{Randall2001}. Excessive fall N carried multi-year legacies: a very high fall rate produced the highest concentrations and losses each year and continued to affect nitrate, residual N, and crop response for more than four years \parencite{Jaynes2015}. Encouragingly, long-term Iowa monitoring showed corn yields rising 144~kg~ha$^{-1}$~yr$^{-1}$ over three decades at a nearly constant rate while drainage nitrate stayed flat, indicating improved N-use efficiency without added input but no reduction in loss from yield gains alone \parencite{Clifford2023}. Modeling reinforces the rate--loss coupling: annual precipitation had the greatest effect on N leaching but fertilizer rate increased all N pools \parencite{Li2006}, and yield plateaued while nitrate loss accelerated quadratically as spring N rose from 100 to 200~kg~ha$^{-1}$ \parencite{Malone2007b}. A cropping-system modeling study found the leaching breakpoint occurred above the agronomic optimum N rate in both continuous maize and maize--soybean, with a larger buffer in the rotation \parencite{Pasley2021}.

\subsection{Nitrogen Timing, Placement, and Enhanced-Efficiency Products}

Shifting N from fall to spring or in-season application generally lowers drainage nitrate, though effects are smaller and less consistent than rate or cropping-system effects. In Minnesota corn--soybean rotations, fall-applied N without nitrapyrin consistently produced the highest concentrations and losses; flow-normalized losses ranked fall $>$ split $>$ spring $=$ fall+nitrapyrin, with reductions relative to fall of 18\% (fall+NP), 17\% (spring), and 13\% (split) in one study \parencite{Randall2003} and 10\%, 14\%, and 6\% (spring+NP) in a companion study \parencite{Randall2005}. Fall N applications produced significantly greater tile nitrate concentration and load than spring applications under corn in Illinois \parencite{Gentry2024}, and spring or side-dress N reduced three-year mean flow-weighted concentrations relative to fall N while maintaining yield \parencite{Pittelkow2017}. Cited Minnesota work found spring application raised N-use efficiency by more than 20\% and reduced annual tile nitrate losses by an average of 36\% relative to fall \parencite{Dinnes2002}. Spring application at the guideline rate produced equal or greater corn yield than a higher fall rate with nitrapyrin \parencite{Vetsch2019}, and fall vs.\ spring effects, while real, were often smaller than effects of weather, drainage volume, crop uptake, and rate \parencite{Liu2021}. Correspondingly, database analyses showed out-of-season and pre-plant timings had higher NO$_3^-$-N concentrations (15.8 and 15.7~mg~L$^{-1}$) than at-plant and side-dress timings \parencite{Hertzberger2019}.

Split and side-dress applications give mixed results. Split fertilization reduced NO$_3^-$-N export at lower rates ($<167$~kg~N~ha$^{-1}$) but not at higher rates \parencite{Liu2021}, and split applications were not significantly different from non-split for annual flow-weighted concentration in a continental synthesis \parencite{Hertzberger2019}. Side-dress timing at V2, V6, or V12 growth stages produced no consistent difference in corn yield or nitrate leaching, with V12 identified as about the earliest reliable sensing stage \parencite{Jaynes2013}. In some cases split N raised loss: a midseason split treatment had consistently higher flow-weighted concentration and greater residual soil nitrate than a single postemergence application \parencite{Jaynes2006}. Notably, at EONR, at-plant applications left less residual N than split (42 vs.\ 62~kg~N~ha$^{-1}$) but showed greater in-season loss, whereas split applications shifted N toward postharvest residual---a nuanced tradeoff \parencite{Bandura2025}. Split-applied urea modestly improved water quality relative to fall and spring anhydrous ammonia (10.0 vs.\ 11.7~mg~L$^{-1}$) without raising yield \parencite{Waring2022}. Variable-rate side-dress N informed by remote sensing lowered observed and simulated concentrations, averaging a 40\% reduction over 15 years attributed mainly to 33\% lower input and timing that avoided peak flow, while maintaining maize yield \parencite{Wilson2020b}. Modeling generally found reducing rate from the recommended level lowered loss with small yield penalties, and that timing affected loss more than yield but without a single optimal timing across sites \parencite{Jeong2018}; watershed simulations found reducing rate 20\% with a 40:60 spring/side-dress split cut nitrate load ~18\% without yield loss \parencite{Singh2023}.

Placement and enhanced-efficiency products offer additional but variable benefits. Fertilizer injection and incorporation showed no significant difference from each other in NO$_3^-$-N export in one synthesis \parencite{Liu2021}, while incorporated and injected applications had higher concentrations than surface applications (14.2, 14.1 vs.\ 10.3~mg~L$^{-1}$) in another \parencite{Hertzberger2019}; injection could increase loss when rate was high \parencite{Jeong2018}. A localized-compaction-and-doming applicator produced 16\% lower flow-weighted concentrations than a spoke injector (16.8 vs.\ 20.1~mg~L$^{-1}$) with similar yield \parencite{Bakhsh2010}. Nitrapyrin generally did not reduce drainage nitrate or improve yield \parencite{Vetsch2019,Pittelkow2017}, though including it with fall N reduced losses 10--18\% in Minnesota rotations \parencite{Randall2003,Randall2005}. Polymer-coated urea increased yield 20\% without drainage but did not consistently outperform non-coated urea for yield \parencite{Nash2015,Nash2015b}. A polymer-coated (ESN) product in Minnesota gave the most robust protection, lowering combined N$_2$O, NH$_3$, and NO$_3^-$ losses 47--49\% below urea and split treatments while maintaining yield, with the largest benefit in wet periods \parencite{Menegaz2025}. Summer fertigation with reduced manure rates cut flow-weighted concentration 16\% versus fall injection (47.7 vs.\ 56.8~mg~L$^{-1}$) and annual load (47.1 vs.\ 55.5~kg~ha$^{-1}$~yr$^{-1}$) via a treatment effect independent of flow, without reducing silage yield \parencite{Gamble2018}. In Missouri, fall anhydrous ammonia with nitrapyrin matched spring anhydrous ammonia for yield and had similar or lower environmental N losses, making it an environmentally acceptable alternative \parencite{Kaur2024,Kaur2024c}. Watershed-scale DSSAT simulations showed the benefit of switching fall-to-spring N for N-use efficiency was largest when cumulative winter rainfall exceeded 500~mm \parencite{Banger2018}.

\subsection{Fertilizer and Manure Source}

Nitrogen source affects drainage nitrate, with manure generally increasing loss relative to inorganic fertilizer at comparable rates. Under continuous corn in Iowa, liquid swine manure produced significantly greater flow-weighted concentrations and losses than UAN with no yield difference, leading authors to recommend corn--soybean rotation for manure management \parencite{Bakhsh2005}. Fall-applied manure with chisel plow gave the highest concentrations, while spring manure with no-till gave lower concentrations, though fall manure also produced the greatest yields \parencite{Bakhsh2007}. In Canada, liquid-hog manure produced higher soil and drainage nitrate than inorganic fertilizer, with more samples exceeding 10~mg~L$^{-1}$, although yields were equivalent \parencite{Elmi2005}; a carrot study found only a small, nonsignificant manure increase and equivalent yield and quality \parencite{Gordon2005}. In Indiana, annually repeated fall-applied liquid swine manure produced the highest nitrate losses among cropped treatments and was the main practice increasing subsurface losses, with time of manure addition controlling loss more than UAN timing/rate or rotation \parencite{HernandezRamirez2011}. Applying swine manure before both corn and soybean, rather than only before corn, raised concentrations and loads \parencite{Dougherty2020}. Fall manure in chisel-plow plots caused more leaching than spring manure in no-till plots \parencite{Rekha2011}. In syntheses, organic fertilizers boosted yields but increased NO$_3^-$-N losses relative to inorganic sources at comparable rates \parencite{Liu2021}, though one meta-analysis found organic sources raised yield without significantly greater dissolved N load \parencite{Christianson2015b}. Among sources, anhydrous ammonia had the highest maximum N$_2$O loss, ammonium nitrate the highest nitrate loss, and urea the highest NH$_3$ loss \parencite{Chatterjee2020b}. Modeling predicted manure gave lower leaching than UAN and spring lower than fall, with fall manure to both crops the highest long-term loss (37.2~kg~N~ha$^{-1}$) \parencite{Malone2007,Malone2007b}.

\subsection{Cover Crops}

Cover crops---particularly winter cereal rye---are among the most consistent and largest in-field reducers of drainage nitrate. In continuous corn in Illinois, split N plus a cereal rye cover crop reduced flow-weighted concentration and load by 35\% and 37\% relative to pre-plant N, and yield-scaled losses by about 33\%, without yield penalty \parencite{PrezaFontes2021}. A companion study found in-season N plus rye decreased leaching 37\% and, on an environmental-damage-cost basis, matched pre-plant N despite shifting losses toward N$_2$O \parencite{PrezaFontes2023}. Cover crops proved more consistent and effective than shifting fertilizer timing: added to spring and fall systems, rye cut flow-weighted concentration 39\% and 38\% and load 40\% and 47\%, reaching 61--64\% reductions in a wet year, and averaged 66\% lower loading than a zero-N control \parencite{Ruffatti2019}. A no-till rye system reduced nitrate concentrations and lowered losses 59\% and 67\% relative to no-till with oat or no-till alone, and altered the load--drainage slope, without reducing yield \parencite{OBrien2022}. Rye after corn significantly reduced mean annual concentration and load under soybean, contributing to a 43\% decrease in cumulative four-year load (from 115 to 65~kg~ha$^{-1}$) \parencite{Gentry2024}. Continuous corn with winter rye had a four-year mean of 5.6~mg~L$^{-1}$, lower than rotation systems, with losses reduced in years of low-to-normal precipitation and high cover-crop growth \parencite{Daigh2015b}. Split N plus rye also gave the lowest dissolved N$_2$O and nitrate loads among fertilized treatments, cutting cumulative N load 56\% versus pre-plant N \parencite{Oliveira2025}. In two Indiana watersheds, cover crops reduced field-scale NO$_3^-$-N losses 58--99\% in most years, pooling to a 65\% reduction, and reduced NH$_4^+$-N losses 45\% \parencite{Vincent2025}.

Cover-crop effectiveness depends on biomass and growing conditions. Biomass and N uptake were low under poor conditions, weakening reductions \parencite{PrezaFontes2021}; an oat cover crop had limited, inconsistent effects, reducing concentration only in a high-biomass year \parencite{OBrien2022}. Cereal rye biomass varied 1.37--5.94~Mg~ha$^{-1}$ with biomass N of 23.5--45.7~kg~ha$^{-1}$, highest under late termination \parencite{Gentry2025}, and growth improved when seeding shifted from post-harvest drilling to broadcasting into standing crop \parencite{Waring2024}. Winter triticale captured 44--164~kg~N~ha$^{-1}$ and reduced profile soil nitrate 33--53~kg~ha$^{-1}$, with 60--100\% of capture occurring by mid-May and minimal residual N below 60~cm \parencite{Nance2007}. Rye maintained higher soil water storage than conventional systems and raised evapotranspiration without a yield disadvantage \parencite{Qi2011}. Some studies noted modest corn yield reductions under rye \parencite{Ruffatti2019,Dougherty2020}, but many report no penalty \parencite{OBrien2022,PrezaFontes2021,Gentry2025}. Modeling estimated a winter cover crop reduced tile nitrate yield 34\% at watershed scale \parencite{David2015}, cut predicted loss 31\% under pre-plant UAN without yield loss \parencite{Malone2007b}, and combined with a 20\% rate cut and split N reduced load 32\% \parencite{Singh2023}; a winter rye cover crop could reduce Chesapeake N loading 11.2~kg~ha$^{-1}$~yr$^{-1}$ \parencite{Woodbury2018}. Under future climate, cover-crop reductions in nitrate loss were projected to persist but may not fully keep pace with rising baseline loss, so cover cropping alone may only partially mitigate climate impacts \parencite{Gupta2023}. At watershed scale, cover crops on 35.5\% of cropland cut subbasin nitrate load 15.5\% and raised median N-use efficiency from 0.49 to 1.02, targeting leaching-dominated areas \parencite{Makhtoumi2025}.

\subsection{Tillage and Residue Management}

Tillage exerts weak and inconsistent control over drainage nitrate loss compared with precipitation, cropping system, and N rate. Over an 11-year Minnesota study, drain flow was higher under no-till than conventional tillage, but nitrate losses were only slightly higher under conventional tillage, and higher no-till flow did not translate into higher nitrate flux \parencite{RANDALL1995}. An 11-year average showed moldboard plow and no-till had similar losses (43 and 41~kg~ha$^{-1}$), and cropping system affected loss far more than tillage \parencite{Randall2001}. In Iowa, tillage only slightly influenced soil nitrate, no-till and ridge-till usually had the lowest drainage concentrations and losses, and crop rotation had greater influence than tillage \parencite{Weed1996}. Tillage had little influence on nitrate losses within a corn--soybean rotation, and continuous corn had higher losses regardless of tillage \parencite{Kanwar1997}. Regional syntheses found conventional and reduced tillage did not differ significantly in drainage N load \parencite{Zhao2016}, though conventional tillage had greater drain discharge while conservation tillage had the highest NO$_3^-$-N concentration (16.1~mg~L$^{-1}$) \parencite{Hertzberger2019}. Modeling likewise found tillage intensity had negligible effects on yield and N$_2$O, with greater intensity slightly reducing drainage flow and N loss \parencite{Wang2016b}.

Where soil-health benefits of reduced tillage accrue, effects on nitrate remain equivocal. No-till improved soil organic matter, aggregate stability, and water-holding capacity in Michigan, but nitrate loads were lower under no-till only in some water years (9.6 vs.\ 18.3~kg~ha$^{-1}$ in one year, higher in another), providing no affirmative evidence that improved soil health reduced leaching in all years, and yields were lower under no-till \parencite{Masri2024}. No-till increased soil organic carbon (0.12~Mg~C~ha$^{-1}$~yr$^{-1}$), reduced the economic optimum N rate by 28~kg~N~ha$^{-1}$, and raised returns \$172~ha$^{-1}$ \parencite{Cassman2025}. Conservation treatments raised soil organic matter and aggregate stability over the long term, and cover crops reduced concentration in both tillage systems (e.g., no-till 8.0 vs.\ 4.4~mg~L$^{-1}$ with rye) \parencite{Waring2024}. Residue and stover management had limited effect on nitrate: residue removal had little impact on concentrations or yields \parencite{Dougherty2020}, 50\% stover harvest reduced simulated drainage N loss only 1.1~kg~ha$^{-1}$ while post-emergence N timing reduced it far more \parencite{Malone2019}, and 38\% stover removal did not significantly change water quality but raised sediment and nutrient losses, offsettable with bioenergy crops \parencite{Guo2018}. Continuous corn with residue removal plus a rye cover crop achieved similar or lower nitrate losses than other row crops \parencite{Daigh2015b}. Reduced tillage had little effect on nitrate but increased dissolved-P loss in one Iranian study \parencite{Farahani2020}.

\subsection{Crop Rotation, Diversification, and Perennial Systems}

Cropping system exerts strong control on drainage nitrate, generally exceeding tillage and often rivaling N rate. Continuous corn consistently lost more nitrate than corn--soybean or perennial systems: continuous-corn plots had much larger losses than rotation plots \parencite{Weed1996}, and continuous corn had higher concentrations and losses than corn--soybean, leading authors to judge continuous corn environmentally unsustainable \parencite{Kanwar1997}. Maize--soybean rotations produced greater yield and lower yield-scaled N loss than continuous maize \parencite{Zhao2016}, and rotation shifted the leaching response with a larger buffer above the agronomic optimum \parencite{Pasley2021}. Changing a southeastern Indiana site from continuous corn to corn--soybean, reducing rate, and adding a cover crop lowered flow-weighted concentration from 28 to 8~mg~L$^{-1}$ over 14 years \parencite{Ritter2015}. Rotation changes reduced simulated tile nitrate loads, more on silt loam than sandy loam \parencite{Ahmed2007b}. More diverse rotations delivered large reductions with maintained yield: a three-year rotation (corn--soybean--wheat/double-crop soybean) had a flow-weighted mean of 4.2~mg~L$^{-1}$ versus 8.0 for a control, with average load 4.9 versus 12.5~kg~ha$^{-1}$, and exceeded 10~mg~L$^{-1}$ in only 1\% of samples---demonstrating sustainable intensification \parencite{Gentry2025}.

Perennial and prairie systems are the most effective at limiting drainage nitrate. Prairie systems had concentrations $\le$1.4~mg~L$^{-1}$, 2--195 times below row crops, and yields 1--11\% of row-crop losses \parencite{Daigh2015b}. Restored prairie grass consistently reduced losses with concentrations near rainfall levels \parencite{HernandezRamirez2011}, and perennial alfalfa and CRP plots had 0.3--4~mg~L$^{-1}$ versus 14--40 for row crops, with four-year Iowa losses of 217, 204, and 202~kg~ha$^{-1}$ for corn systems versus 7 and 4 for alfalfa and CRP \parencite{Randall2001}. Perennial plots had the lowest concentration and loading among treatments \parencite{Waring2024}. However, the benefit of perennials can be short-lived on conversion back to row crops: nitrate losses were low only in the first year, then similar to row crops, as residual soil nitrate rose sharply, though former-alfalfa treatments averaged 20\% less drain flow \parencite{Huggins2001}. Simulated novel perennial systems (intermediate wheatgrass, alfalfa) cut tile nitrate loads ~91--92\% in perennial years, and converting all corn--soybean area to alfalfa or wheatgrass reduced watershed losses 78--81\% \parencite{Wilson2023}. Bioenergy perennials also reduce loss: strategically implemented miscanthus cut total N leaching and DIN export 19--26\% under future climate \parencite{Ferin2023}; replacing maize with switchgrass could reduce Chesapeake N loading 18.2--18.9~kg~ha$^{-1}$~yr$^{-1}$ \parencite{Woodbury2018}; and bioenergy crops on marginal land slightly reduced watershed nutrient losses \parencite{Guo2018}. Diversified rotations, cover crops, perennials, wetlands, bioreactors, and drainage control were all recognized as complementary reducers \parencite{Dinnes2002}. Drainage often raised soybean and corn yield modestly (e.g., +6\% soybean, +8\% corn without N) \parencite{Cassman2025}, though on Fargo clay drainage sometimes reduced yield in dry years, and corn--soybean consistently outyielded continuous corn \parencite{Acharya2019b}. Rotation raised corn yield without changing leaching over six years \parencite{Bakhsh2005}, and organic rotations offered water-timing benefits with sometimes comparable yields \parencite{Logsdon2021}.

\subsection{Drainage Water Management and Controlled Drainage}

Managing the drainage system itself---controlled drainage, subirrigation, and water-table control---reduces drainage volume and nitrate export, though with pathway tradeoffs. Simulated drainage water management reduced annual drain flow ~56\% and drain-flow nitrate loss ~46\% for both continuous corn and corn--soybean, but increased nitrate in surface runoff 3--4~kg~ha$^{-1}$, so net reduction reaching surface water was smaller (32--38\% when combined with reduced rate) \parencite{Ale2012b}. A subirrigation/controlled-drainage system in Ohio lowered deep-profile soil nitrate 46\%, raised corn and soybean yields 19\% and 64\%, and increased N uptake, reducing nitrate available for export \parencite{Fisher1999}. In Missouri, free and managed drainage raised corn yield 15--21\% over no drainage, with managed drainage averaging a further 4\% \parencite{Nash2015}, though managed drainage did not consistently outyield free drainage across years \parencite{Nash2015b}. In Chinese paddy fields, raising the water-table control level in the first week after fertilization reduced total N losses ~9.5\%, mainly via reduced subsurface outflow, making combined controlled irrigation and controlled drainage effective \parencite{He2019b}. Controlled drainage slightly increased corn yield and reduced drainage flow, N loss, and concentration, but increased N$_2$O \parencite{Wang2016b}. Controlled drainage was repeatedly identified as a supplementary mitigation option where in-field management is insufficient \parencite{HernandezRamirez2011}. A caveat: on an existing field system, drainage water management did not reduce load because water moved laterally to an adjacent free-draining system \parencite{David2015}. In arid maize systems, demand-based irrigation with intermittent leaching (threshold ~11~g~kg$^{-1}$ in the 60--100~cm layer) controlled salinity and sustained yield, whereas insufficient leaching cut yield ~86.8\% over 12 years \parencite{Liu2024}.

\subsection{Irrigation and Vegetable Systems}

In irrigated and specialty systems, water application governs nitrate leaching. Sprinkler irrigation produced lower total nitrate and dissolved P than flood and furrow irrigation and, combined with biochar, reduced nitrate concentration, making sprinkler an appropriate water-conserving practice \parencite{Farahani2020}. Greater absolute nitrate loss under furrow than sprinkler reflected lower furrow efficiency \parencite{Hurtado2011}. In drip-irrigated mulched vegetables in Florida, increasing irrigation volume and N above recommendations raised nitrate leaching 32--66\% across crops, and drainage lysimeters and soil cores agreed while ceramic suction cups underestimated leaching, especially under wetter treatments \parencite{Zotarelli2007}. Soil amendments had localized effects: phosphogypsum decreased and refuse lime increased flow-weighted concentrations in some drainage events, without significant differences in total seasonal loss \parencite{Hurtado2011}.

\subsection{Hydrologic and Climatic Controls}

Hydrology---drainage amount and its seasonal timing, driven by precipitation---is the dominant control on interannual and among-plot variability in nitrate loss. Flow-weighted concentrations rose sharply in wet years following dry periods, and losses were lowest in dry years because of low drain flow; the tillage effect was small relative to precipitation \parencite{RANDALL1995}. Minnesota drainage averaged 38~mm during a dry period but 507~mm during a wet period; corn--soybean plots ranged from $<$50~mm and 7--18~mg~L$^{-1}$ in dry years to $>$350~mm and ~35~mg~L$^{-1}$ in a wet year \parencite{Randall2001}. Loss is strongly seasonal: 62\% of annual drainage and ~68--70\% of annual nitrate loss occurred in April--June in Minnesota \parencite{Randall2003}, with 71\% of drainage and 73--77\% of loss in April--June in a companion study \parencite{Randall2005}. In Iowa, ~60\% of tile flow and nitrate leaching occurred March--May \parencite{Bakhsh2007}, and the March--June share of seasonal nitrate loss ranged 25--99\% across years \parencite{Bakhsh2010}. Spatiotemporal variability produced three- to fourfold differences in flow and nitrate loss even under identical management \parencite{Bakhsh2010}. Concentrations were significantly affected by seasonal rainfall and crop conditions \parencite{Bakhsh2005}, nitrate load and drainage were tightly correlated (0.86--0.99) while precipitation correlations were weaker \parencite{Masri2024}, and nitrate loads were linearly correlated with cumulative annual drainage, with management altering the slope \parencite{OBrien2022}. Regionally, relative precipitation had a strong exponential effect: at 148\% of mean precipitation, drainage N loss reached 67.5~kg~ha$^{-1}$ \parencite{Zhao2016}, and wetter conditions were associated with greater dissolved N losses, though very high rates could overshadow precipitation \parencite{Christianson2015b}. Even with no N applied, nitrate loss remained substantial (~8.3~mg~L$^{-1}$), reflecting background soil N mineralization \parencite{Waring2022}, and non-growing-season mineralization from soil and residues is an important nitrate source \parencite{Gentry2024}. Most drainage and nitrate loss occurred outside the growing season in cool humid systems \parencite{Gordon2005}. Long-term monitoring is essential: short 5--10-year windows can show apparent trends differing from the full record \parencite{Clifford2023}, and tile systems respond within 6--14 months but excessive fall N legacies persist for years \parencite{Jaynes2015}. Weather also modulates enhanced-efficiency products, with ESN lowering N$_2$O chiefly in wetter periods \parencite{Menegaz2025}.

\subsection{Watershed-Scale Programs, Soil Testing, and Nitrogen-Use Efficiency}

Watershed-scale implementation of soil-test-based N programs and prioritized conservation confirms in-field mechanisms at landscape scale. Adopting the late-spring nitrate test (LSNT) across a 400-ha Iowa subbasin reduced drainage nitrate concentration by at least 30\%, associated with reduced fertilizer use, while corn yields matched non-limiting strips \parencite{Jaynes2004}. LSNT at watershed scale produced significant crop responses, with stalk nitrate and SPAD useful indicators, and was effective for reducing loss though it may raise producer risk \parencite{Karlen2005}. Modeling reproduced the LSNT-driven concentration decrease, giving long-term losses of 20.4, 22.2, and 27.3~kg~N~ha$^{-1}$ for LSNT, single spring, and single fall applications, with early-season weather explaining 90\% of interannual LSNT-rate variation \parencite{Malone2010}. Soil-test programs (PSNT, LSNT) generally reduced residual soil nitrate and loss versus yield-goal management \parencite{Dinnes2002}, and changing timing and reducing rate through soil testing cut concentrations 30\% in a central Iowa subbasin \parencite{Ritter2015}. Field-level prioritization tools scaled cover crops and grassland conversion, raising watershed N-use efficiency and cutting nitrate load, with conservation tillage least effective \parencite{Makhtoumi2025}. Alternative farming practices reduced drainage 41\% and concentration from 17.2 to 8.2~mg~L$^{-1}$, cutting nitrate losses 59--62\% and staying below 10~mg~L$^{-1}$ \parencite{Oquist2007}. Because 85\% of a Minnesota watershed's nitrate load came from tile drainage, substantial perennialization was needed for a 30--50\% watershed reduction \parencite{Wilson2023}. Adoption remains a constraint: farmers commonly used soil testing but were less familiar with cover crops, controlled drainage, wetlands, and bioreactors, with financial, operational, and social barriers \parencite{David2015}. N-use efficiency and partial factor productivity were highest for side-dress and lowest for very high fall N \parencite{Jaynes2015}, and EONR at an N:corn price ratio of 5.6 is proposed as a sustainability standard balancing profit and environmental risk \parencite{Bandura2025}.

\subsection{Process Understanding, Simulation Models, and Loss-Pathway Tradeoffs}

Plant uptake mechanistically underlies most management effects: plots with plants had lower soil and vadose-water nitrate than barren plots at 30--90~cm, confirming that living cover removes nitrate \parencite{Bosompemaa2021}. Process-based models---RZWQM, DRAINMOD-N II, DNDC, DSSAT, Daisy, APSIM, SWAT, CoupModel---reproduce tile flow, nitrate loss, and yield with satisfactory-to-excellent skill and are used to extrapolate management scenarios \parencite{Jeong2018,Ahmed2007b,Ale2012b,Pasley2021,Motarjemi2021,Conrad2009}. A modified DNDC captured discharge and leaching magnitudes across fertilizer treatments \parencite{Li2006}; Daisy simulated cumulative discharge and nitrate leaching well (NSE 1.00 and 0.87) though daily fluxes and frost/thaw periods were harder \parencite{Motarjemi2021}; DRAINMOD-N II reproduced monthly losses and quantified that ~20\% rate reduction cut losses 10--17\% \parencite{Ale2012b}; and regression equations explained 76--95\% of variation in yield, drainage, concentration, and loss \parencite{Malone2007}. RZWQM applications identified reduced rate and spring/split timing as effective, with 150~kg~ha$^{-1}$ single spring N favorable \parencite{Azevedo1997,Malone2007b}. Climate-adaptation modeling projected declining yield and rising drainage N loss under warming, with an optimum rate near 120~kg~N~ha$^{-1}$ balancing economics and environment, and identified cultivar, planting date, tillage, and controlled drainage effects \parencite{Wang2016b}.

Finally, mitigation practices redistribute N among loss pathways, so single-pathway evaluation can mislead. In-season N plus rye cut leaching 37\% but raised annual N$_2$O 26\%, leaving total N loss and environmental-damage cost (519--612~US\$~ha$^{-1}$ among fertilized treatments) statistically unchanged versus pre-plant N \parencite{PrezaFontes2023}. Split N plus rye had the lowest dissolved N$_2$O and nitrate loads, showing stacked practices need not trade one pathway for another \parencite{Oliveira2025}, and no consistent nitrate--N$_2$O tradeoff emerged across tillage and cover-crop systems \parencite{OBrien2022}. ESN protected against N$_2$O, NH$_3$, and NO$_3^-$ simultaneously (47--49\% lower combined loss) while maintaining productivity \parencite{Menegaz2025}. N$_2$O emissions rose nonlinearly with N rate, with a threshold near 120~kg~N~ha$^{-1}$ \parencite{Wang2016b}, and each unit of fertilizer-N produced 0.01~kg~N$_2$O-N and 0.12~kg~NO$_3^-$-N loss, with subsurface-drained fields having the highest mean N$_2$O loss \parencite{Chatterjee2020b}. Emerging monitoring links crop performance to losses: early-season NDVI was negatively correlated with cumulative N$_2$O, suggesting remote sensing can flag poorly performing fields \parencite{Fontes2019}. Economic framings increasingly accompany these tradeoffs, from environmental-damage-cost accounting \parencite{PrezaFontes2023} to soil-health return analyses valuing no-till and drainage \parencite{Cassman2025} and switchgrass reference prices for N reduction \parencite{Woodbury2018}. The consistent conclusion across scales is that combining two or more practices---rate optimization, in-season timing, cover crops or diversified rotation, and drainage management---works in harmony to reduce nitrate loss most cost-effectively, since in-field management alone is often insufficient \parencite{Dinnes2002,Singh2023}.